% COL764 A4 report - at most 2 pages. Compile: pdflatex report.tex  (or upload to Overleaf), then copy
% report.pdf to the repository root so that make_submission.py picks it up.
% Every TODO must be replaced with numbers from your own runs (dev/compare.py, dev/analyze.py).
\documentclass[10pt,a4paper]{article}
\usepackage[margin=1.8cm]{geometry}
\usepackage{booktabs,microtype,enumitem}
\setlist{nosep,leftmargin=*}
\setlength{\parskip}{3pt}\setlength{\parindent}{0pt}

\title{\vspace{-1.6cm}COL764 Assignment 4: Iterative Retrieval for Multi-hop QA}
\author{TODO Name \quad TODO Entry number}
\date{\vspace{-0.8cm}}

\begin{document}
\maketitle

\section{Single-shot vs.\ iterative (dev set, TODO $n$ questions)}
\begin{center}
\begin{tabular}{lcccccc}
\toprule
System & EM & Answer F1 & Evidence F1 & Searches / q & LLM calls / q & $S$ \\
\midrule
% paste the LaTeX rows printed by: python dev/compare.py runs/<b1> runs/<ours> --names B1 Ours
B1 (single-shot) & TODO & TODO & TODO & 1.00 & 1.00 & TODO \\
Ours (iterative) & TODO & TODO & TODO & TODO & TODO & TODO \\
\bottomrule
\end{tabular}
\end{center}
Retrieval diagnostics (from \texttt{dev/analyze.py}): both gold paragraphs are retrieved for TODO\% of
questions (B1: TODO\%), and the agent's bridge entity matches the gold bridge for TODO\%.

\section{Design}
\textbf{Plan and \texttt{next\_query}.} One LLM call decomposes the question into two single-hop
sub-questions, the second referring to the first answer as \#1 (four in-context examples). Hop~1 searches
both $Q_1$ ($k{=}6$) and the original question ($k{=}4$): the latter's extra context words often rank the
bridge paragraph when $Q_1$ alone is too short for BM25. After hop~1 is read, the bridge entity is
substituted into $Q_2$ and the agent issues two queries: the filled-in $Q_2$ ($k{=}6$) and the bridge
entity on its own ($k{=}4$), which retrieves the bridge's own article, where hop-2 facts (spouse,
birthplace, founding date, \dots) usually live. If a read fails, the next query is the current
sub-question prefixed with the title of the best retrieved passage not yet used (a candidate bridge
entity) at $k{=}10$, so queries always depend on earlier results and never repeat. Searches are cheap
and LLM calls scarce, so the budget is spent on two complementary queries per hop rather than more reads.

\textbf{\texttt{should\_stop} and \texttt{answer}.} Each round ends with one \emph{read} call that
asks for the supporting passage number first and then the shortest answer span (evidence-first output
grounds the 4B model). The first read of a hop may say ``unknown''; that triggers another round, whose
read must give a best guess. The agent stops as soon as hop~2 is answered, or when only one LLM call or
less than 15\,s remains; that last call is reserved for a forced best-guess answer over all evidence. The
two citations are the passages that the hop-1 and hop-2 reads point to (each checked to actually contain
its answer, else the best-ranked passage that does): gold evidence is exactly two paragraphs, so citing
two rather than three raises evidence precision. A typical question uses TODO searches and TODO LLM calls.

\textbf{What did not help.} TODO (e.g.\ a free-form chain-of-thought read: longer outputs, slower, and
answer F1 dropped from TODO to TODO because the model wrapped answers in sentences).

\section{Error analysis}
Ten dev questions the agent gets wrong (EM $=0$), labelled by hand (starting from the automatic labels of
\texttt{dev/analyze.py}):
\begin{center}
\begin{tabular}{lc}
\toprule
Category & Count \\
\midrule
(i) bridge entity not found & TODO \\
(ii) stopped too early & TODO \\
(iii) evidence retrieved but answered wrongly & TODO \\
(iv) answer correct but marked wrong & TODO \\
(v) other & TODO \\
\bottomrule
\end{tabular}
\end{center}
\textbf{Example 1 (TODO category).} TODO question id, question, gold vs.\ predicted answer, what the
queries retrieved, and why it failed.

\textbf{Example 2 (TODO category).} TODO.

\section{AI-use disclosure and code provenance}
TODO: state honestly which tools you used and for what (e.g.\ an AI coding assistant helped draft
\texttt{agent.py} and the analysis scripts in \texttt{dev/}; all design choices, experiments, numbers and
this report were checked by me). The starter code (\texttt{a4/}, \texttt{run\_dev.py}, the baseline)
is from the course repository.

\end{document}
